\begin{abstract}
We characterize the minimax mean-squared error for stationary policy evaluation from one stationary behavior trajectory in a partially observed Markov decision process with binary observed states, binary hidden states, and binary actions. The behavior and target policies are known, actions are randomized conditional on the observed state, rewards lie in \([0,1]\), action likelihood ratios are bounded, and both policy-induced joint-state kernels satisfy a common strict total-variation contraction bound. At fixed mixing and action-overlap parameters, the minimax risk has order \(T^{-1}\), where \(T\) is the observed trajectory length. Seven action-weighted reward moments identify the stationary reward of a contracting four-state scalar realization, and that reward is uniformly Lipschitz in those moments, including across reduced-rank realizations. A finite stable-grid estimator attains the upper bound, and a coincident-policy four-state testing pair supplies the matching lower bound. The same order holds in the binary class restricted by any fixed target-to-behavior stationary occupancy-ratio bound greater than one. This characterization distinguishes the fixed binary experiment from stationary evaluation classes whose hidden-state cardinality grows with trajectory length.
\end{abstract}

\section{Introduction}

This paper establishes a matched parametric minimax rate for stationary policy evaluation in a fixed binary partially observed experiment. The data consist of one stationary behavior trajectory recording a binary observed state, a binary action, and a reward in \([0,1]\) at each of \(T\) periods, where \(T\) is the observed trajectory length. A binary hidden state completes the four-point joint-state space. The behavior and target policies are supplied, and behavior actions are randomized through the observed state conditional on the full pre-action history. Under bounded action likelihood ratios and a common strict total-variation contraction bound for the behavior and target joint-state kernels, the minimax mean-squared error for the stationary target mean reward has order \(T^{-1}\).

Stationary policy evaluation connects sequential causal interventions to long-run outcomes in controlled stochastic systems \citep{Robins1986,Murphy2003,Puterman1994}. Hidden states allow rewards and transitions to depend on information beyond the recorded state. In the experiment studied here, observed-state randomization makes short action-weighted reward histories informative about target-policy interventions. Fixed joint-state dimension then supplies a finite recurrence that connects those intervention moments to the stationary reward. The resulting statistical guarantee concerns the scalar policy value uniformly over the binary class in \cref{obj:def:binary-pomdp-class}.

The rate characterization holds at each fixed positive mixing scale \(t_0\), the parameter specifying the contraction bound, and fixed nonnegative logarithmic action likelihood-ratio bound \(\zeta\). The corresponding bounds are \(\alpha=\exp(-1/t_0)\), the contraction factor, and \(L=\exp(\zeta)\), the action likelihood-ratio bound. For every \(T\geq12\), \cref{obj:thm:fixed-binary-minimax} bounds the observed-data minimax risk below by \(1/(1024T)\) and above by a constant depending on \(\alpha\) and \(L\), divided by \(T\). The joint reward and successor-state law may couple the current reward with the next state. The guarantee also covers reduced-rank realizations under the stated contraction conditions.

The construction uses seven action-weighted reward moments. Their expectations follow successive target-policy transitions, while contraction turns the resulting four-state recurrence into a uniform stationary-value modulus, including at reduced-rank realizations. The stable-grid estimator in \cref{obj:def:stable-grid-estimator} fits these moments and attains the risk bound in \cref{obj:thm:stable-grid-upper}; a four-state local testing pair supplies the matching converse. The displayed estimator has a finite candidate list of size \(O(T^{19})\) at fixed mixing scale, as quantified by \cref{obj:prop:exhaustive-grid-arithmetic-complexity}; this is a candidate-count statement, not a runtime bound.

The closest stationary-trajectory comparison is \citet[Theorem 2.1, Corollary 2.2, and Theorem 3.1]{HuWager2023}. Their partial-history importance-weighting analysis balances geometric forgetting against the variance cost of longer weighting windows. At fixed mixing and action-overlap parameters, their single-trajectory minimax rate has exponent \(2/(2+t_0\zeta)\), over an experiment permitting hidden-state cardinality to increase with trajectory length. The public manuscript \citet[Theorem 1 (Uniform overlap frontier) and the Signed-depth family definition]{CausalSmithLatentOverlap2026} supplies a cited cardinality-uniform comparison with bounded stationary occupancy ratios; its hard family likewise uses increasing hidden depth. For positive action-overlap exponent, the fixed binary characterization identifies the distinct rate obtained when the hidden alphabet contains two states throughout.

The information conditions also distinguish this experiment from proxy-based evaluation. \citet[Section 4, Assumptions 2--4, and Theorems 2--4]{NairJiang2021} use observable matrices with rank equal to hidden-state cardinality for finite-horizon identification. \citet[Section 4 and Theorems 7--8]{ShiUeharaHuangJiang2022} develop discounted-value bridge estimation using independent observable transition and proxy tuples. Here, observed-state randomization identifies intervention moments from one dependent trajectory, and contraction stabilizes their stationary scalar limit. The finite-moment mechanism connects to classical realization theory \citep[pp.\ 545--548]{Ho1966}\citep[pp.\ 26--40]{Carlyle1971}.

The fixed-binary risk bounds and moment identities are proved under the stated law conditions; the cardinality-uniform comparison uses a cited public-manuscript conclusion.

The paper proceeds through related work, setup and assumptions, and the main results. The discussion develops the statistical and candidate-count implications, followed by limitations and future work. The appendices present the supporting mathematical statements and proofs of the main results.
